% 复用已校验的题解 Markdown。
% @card {"id":"td2-1-1","type":"exercise","title":"1.1 · 分割生成的 σ-代数","fr":"Tribu engendrée par une partition","refs":[["td2",2,2],["answer2",2,2]],"requires":["partition-sigma","generated-sigma"],"import_solution":"1.1"}
% @end
% @card {"id":"td2-1-2","type":"exercise","title":"1.2 · 自然数上的原子测度","fr":"Une mesure atomique sur N","refs":[["td2",2,2],["answer2",3,4]],"requires":["measure-definition","point-atoms","finite-sigma-finite"],"import_solution":"1.2"}
% @end
% @card {"id":"td2-1-3","type":"exercise","title":"1.3 · 代数不一定是 σ-代数","fr":"Une algèbre qui n’est pas une tribu","refs":[["td2",2,2],["answer2",5,5]],"requires":["finite-cofinite","sigma-algebra"],"import_solution":"1.3"}
% @end
% @card {"id":"td2-1-4","type":"exercise","title":"1.4 · 从有限可加到可数可加","fr":"Caractérisation de la σ-additivité","refs":[["td2",2,2],["answer2",6,6]],"requires":["increasing-continuity","measure-definition"],"import_solution":"1.4"}
% @end
% @card {"id":"td2-1-5","type":"exercise","title":"1.5 · 计数测度","fr":"La mesure de comptage","refs":[["td2",2,2],["answer2",7,7]],"requires":["measure-definition","subadditivity"],"import_solution":"1.5"}
% @end
